\documentclass[../main.tex]{subfiles}
\begin{document}
%==========================================TIÊU ĐỀ TẬP========================================
\begin{table}[h]
  \begin{adjustwidth}{-1cm}{}
    \begin{tabular}{>{\centering\arraybackslash}p{6cm}|>{\raggedright\arraybackslash}p{12cm}}
      \multirow{5}{6cm}
      % Cột bên trái: ÉPISODE và số tập
      {\\[-40pt]
        \centering
        {\fontsize{20pt}{20pt}\selectfont{\outlinetext[1]{eptype}{white}{\flaregothic ÉPISODE}}\\[12pt]
          {\fontsize{72pt}{72pt}\selectfont{\outlinetext[1]{epnum}{white}{\burbank 157}}}\\[6pt]       % use for normal/special episodes
          % {\fontsize{36pt}{36pt}\selectfont{\outlinetext[1]{epnum}{white}{\burbank SÁU}}}\\[6pt]     % use for custom episodes
          {\fontsize{20pt}{20pt}\selectfont{\outlinetext[1]{epnum}{white}{\cafeteria (12-8)}}}\color{black}}
        \\[18pt]
        % {\fontsize{14pt}{14pt}\selectfont{\outlinetext[1]{epattr}{white}{\flaregothic [I]}}}\\
        {\fontsize{18pt}{18pt}\selectfont{\outlinetext[1]{epattr}{white}{\flaregothic ÉPISODE PERDU}}}\\
        \color{black}
      }
       &
      % Dòng thứ nhất: tên tập tiếng Nhật
      {\fotskip\fontsize{18pt}{18pt}\selectfont \ruby{全}{ぜん}\ruby{部}{ぶ}お\ruby{前}{まえ}のせいだろ！}                  \\[6pt]
      % Dòng thứ hai: tên tập tiếng Anh
       & {\cafeteria\fontsize{18pt}{18pt}\selectfont All Your Fault!}              \\[4pt]
      % Dòng thứ ba: tên tập tiếng Việt
       & {\vnmsans\fontsize{18pt}{18pt}\selectfont Chiếc nơ nịnh nọt}             \\[8pt]
      % Dòng thứ tư: Nhân vật xuất hiện
       & {\caviar{\fontsize{14pt}{14pt}\selectfont Nhân vật: \emp{kuon2}}} \\
      % Dòng thứ năm (nếu có): Nhân vật chỉ xuất hiện trong Omake
       & {\abv Truyện phụ:} \emp{okayu}
      \\[8pt]
      % Dòng cuối cùng: Thời gian diễn ra của tập (thường sẽ là ngày phát hành)
       & {\sen\fontsize{16pt}{16pt}\selectfont Mốc thời gian: 15/05/2022}  \\[18pt]
    \end{tabular}\\[8pt]
  \end{adjustwidth}
\end{table}
%=========================================TRUYỆN CHÍNH========================================
\color{black}

\txtbx[2]{
  {\color{vol} \uline{Tóm tắt:}} ???
}

%==========================================TRUYỆN PHỤ=========================================
\omk

\end{document}